The notion of parameter-efficient fine-tuning appeared in a series of pioneering works \cite{rebuffi2017, rebuffi2018, rosenfeld2020} that address the problem of transfer learning in the computer vision domain. These authors developed the idea of an adapter, a small learnable auxiliary module, specifically as a small convolution layer plugged in parallel to the original convolution layer.



The adapter concept was extended to the natural language processing (NLP) domain by \cite{houlsby2019} and has since been studied in various works \cite{dettmers2023, he2022, hu2021, pfeiffer2021, ruckle2021, zhang2023}. In \cite{houlsby2019}, an adapter, represented as a bottleneck residual block, is appended at the end of each Transformer sub-block. The adapter’s residual block consists of two fully connected layers with intermediate feature dimensions $m < d$.

Several studies \cite{dettmers2023, hu2021, lialin2023, zhang2023} have focused on low-rank adapters. \cite{hu2021} pioneered this area with LoRA, a simple adapter structure plugged in parallel, which consists of the multiplication of two matrices of incomplete rank. \cite{zhang2023} proposed an adaptive variant of LoRA, called AdaLoRA, which uses Singular Value Decomposition (SVD) for adapter representation and rank truncation. Practical aspects of a low-rank adapter-based approach are studied in \cite{cherniuk2023, dettmers2023}. \cite{cherniuk2023} presented a scheme for training a low-rank adapter with original model weights in quantized arithmetics, while \cite{dettmers2023} observed that a careful choice of matrix multiplication order in low-rank adapters significantly reduces computational and memory usage.

Alternative research directions to adapters have also been explored \cite{benzaken2022, guo2021, lester2021, li2021, liu2022}. \cite{benzaken2022} reported that even a small number of parameters in bias vectors are sufficient to achieve competitive performance. DiffPruning, introduced by \cite{guo2021}, imposes an $L_0$-regularization constraint with differential relaxation on a vector of auxiliary parameters $\delta \tau$ to reparametrize original parameters. \cite{liu2022} proposed scaling intermediate activations in attention and linear layers. The connection between prompt fine-tuning methods \cite{lester2021, li2021, liu2023} and adapter methods is revealed in \cite{he2022}.
